% =========================================================================
% CABEÇALHO UNIFICADO E DESIGN SYSTEM (NETWORKS FLOW)
% =========================================================================

% =========================================================================
% PACOTES FUNDAMENTAIS
% =========================================================================
\usepackage{amssymb}
\usepackage{amsmath}
\usepackage{amsthm}
\usepackage{bm}
\usepackage{enumitem}
\usepackage{mathtools}
\usepackage{lipsum}
\usepackage[brazilian]{babel}
\usepackage[utf8]{inputenc}
\usepackage[T1]{fontenc}
\usepackage[pdftex]{graphicx}
\usepackage{color}
\usepackage{xcolor}
\usepackage{colortbl}
\usepackage{tikz}
\usetikzlibrary{arrows.meta, calc, positioning, decorations.markings}
\usepackage{verbatim}
\usepackage[hyphens]{url}
\usepackage{longtable}
\usepackage{booktabs}
\usepackage[section]{placeins}
\usepackage{listings}
\usepackage{helvet}
\renewcommand{\familydefault}{\sfdefault}
\usepackage{titlesec}
\usepackage{setspace}
\usepackage[plainpages=false,pdfpagelabels]{hyperref}

\hypersetup{
    colorlinks=true,
    linkcolor=blue!80!black,
    citecolor=green!50!black,
    urlcolor=blue!60!black
}

% =========================================================================
% LISTAGENS DE CÓDIGO-FONTE C++ (LISTINGS)
% =========================================================================

% -------------------------------------------------------------------------
% Paleta de Cores de Sintaxe
% -------------------------------------------------------------------------
\definecolor{codegreen}{rgb}{0,0.6,0}
\definecolor{codegray}{rgb}{0.5,0.5,0.5}
\definecolor{codepurple}{rgb}{0.58,0,0.82}
\definecolor{backcolour}{rgb}{0.9,0.9,0.9}

% -------------------------------------------------------------------------
% Configurações Globais do Listings
% -------------------------------------------------------------------------
\lstset{
    language=C++,
    basicstyle=\ttfamily\small,
    tabsize=4,
    keepspaces=true,
    columns=fullflexible,
    lineskip=-2pt,
    backgroundcolor=\color{backcolour},
    commentstyle=\color{codegreen},
    keywordstyle=\color{blue}\bfseries,
    stringstyle=\color{codepurple},
    deletestring={[b]'},
    numbers=left,
    numberstyle=\tiny\color{codegray},
    numbersep=6pt,
    breaklines=true,
    breakatwhitespace=false,
    showspaces=false,
    showstringspaces=false,
    showtabs=false,
    frame=single,
    rulecolor=\color{black},
    captionpos=b
}

% =========================================================================
% ESTRUTURAS SEMÂNTICAS, TEOREMAS E DEFINIÇÕES
% =========================================================================
\newtheorem{prop}{Proposição}

% -------------------------------------------------------------------------
% Definições Formais
% -------------------------------------------------------------------------
\newcounter{definicao}
\renewcommand{\thedefinicao}{\arabic{definicao}}
\NewDocumentEnvironment{definicao}{o o}{%
    \par\addvspace{0.8\baselineskip}\noindent
    \refstepcounter{definicao}%
    \IfValueT{#2}{\label{#2}}%
    \textbf{Definição~\thedefinicao\IfValueT{#1}{ (#1)}}.%
    \space
}{%
    \par\addvspace{0.5\baselineskip}
}

% -------------------------------------------------------------------------
% Teoremas e Demonstrações
% -------------------------------------------------------------------------
\newcounter{teorema}
\renewcommand{\theteorema}{\arabic{teorema}}
\NewDocumentEnvironment{teorema}{o o}{%
    \par\addvspace{0.8\baselineskip}\noindent
    \refstepcounter{teorema}%
    \IfValueT{#2}{\label{#2}}%
    \textbf{Teorema~\theteorema\IfValueT{#1}{ (#1)}}.%
    \space\itshape
}{%
    \par\addvspace{0.5\baselineskip}
}

\NewDocumentEnvironment{demonstracao}{}{%
    \par\noindent\textbf{Demonstração.}\space
}{%
    \hfill $\blacksquare$\par\addvspace{0.5\baselineskip}
}

% -------------------------------------------------------------------------
% Macros Conceituais e Auxiliares
% -------------------------------------------------------------------------
\newcommand{\intuicao}[1]{\paragraph{Intuição #1.}}
\newcommand{\implicacoes}[1]{\paragraph{Implicações Estratégicas #1.}}
\newcommand{\sa}{\text{s.a.}\quad}

% =========================================================================
% PALETA CROMÁTICA SEMÂNTICA (DESIGN SYSTEM)
% =========================================================================

% -------------------------------------------------------------------------
% Partições S e T
% -------------------------------------------------------------------------
\definecolor{setSborder}{RGB}{22, 101, 52}     % [COR] Borda e contorno estrutural de S (Emerald 800)
\definecolor{setStext}{RGB}{20, 83, 45}        % [COR] Tipografia de alto contraste para rótulos de S (Emerald 900)
\definecolor{setSsource}{RGB}{187, 247, 208}   % [COR] Fundo luminoso para vértice fonte s e sementes (Emerald 200)
\definecolor{setSnode}{RGB}{220, 252, 231}     % [COR] Fundo pastel para vértices internos de S (Emerald 100)
\colorlet{setSbg}{green!10}                    % [COR] Fundo pastel suave para região de S / Foreground

\definecolor{setTborder}{RGB}{190, 18, 60}     % [COR] Borda e contorno estrutural de T (Rose 700)
\definecolor{setTtext}{RGB}{159, 18, 57}        % [COR] Tipografia de alto contraste para rótulos de T (Rose 900)
\definecolor{setTsink}{RGB}{254, 205, 211}     % [COR] Fundo luminoso para vértice sumidouro t e sementes (Rose 200)
\definecolor{setTnode}{RGB}{255, 228, 230}     % [COR] Fundo pastel para vértices internos de T (Rose 100)
\colorlet{setTbg}{red!10}                      % [COR] Fundo pastel suave para região de T / Background

% -------------------------------------------------------------------------
% Fronteira de Corte e Elementos Neutros
% -------------------------------------------------------------------------
\definecolor{cutviolet}{RGB}{109, 40, 217}     % [COR] Violeta imperial para fronteira e arestas do corte mínimo
\definecolor{cutred}{RGB}{225, 29, 72}          % [COR] Vermelho rubro padrão para fronteira de corte mínimo
\definecolor{darkslate}{RGB}{30, 41, 59}        % [COR] Slate escuro para tipografia e arestas neutras

% -------------------------------------------------------------------------
% Cores da Capa e Cards
% -------------------------------------------------------------------------
\definecolor{topbg}{RGB}{15, 23, 42}           % [COR] Slate 900 profundo da divisória da capa
\definecolor{bottombg}{RGB}{248, 250, 252}     % [COR] Fundo inferior: Papel Neve (Slate 50)
\definecolor{accentgold}{RGB}{226, 192, 116}   % [COR] Dourado nobre metálico
\definecolor{accentcyan}{RGB}{14, 165, 233}    % [COR] Azul ciano vibrante
\definecolor{cardbg}{RGB}{224, 232, 242}       % [COR] Fundo do card do grafo
\definecolor{cardborder}{RGB}{186, 202, 220}   % [COR] Borda do card do grafo

% -------------------------------------------------------------------------
% Aliases de Compatibilidade para a Paleta Semântica
% -------------------------------------------------------------------------
\colorlet{scyan}{setSborder}
\colorlet{scyanbg}{setSnode}
\colorlet{ssourcebg}{setSsource}
\colorlet{tgold}{setTborder}
\colorlet{tgoldbg}{setTnode}
\colorlet{tsinkbg}{setTsink}

% =========================================================================
% BIBLIOTECA DE ESTILOS GRÁFICOS (TIKZ)
% =========================================================================
\newcommand{\tikzsubcaption}[2]{\node[font=\footnotesize\bfseries, above] at #1 {#2};}

\tikzset{
    >=Stealth,
    % ---------------------------------------------------------------------
    % Redes de Fluxo e Otimização
    % ---------------------------------------------------------------------
    flow node/.style={circle, draw, fill=white, minimum size=8.5mm, inner sep=1pt, font=\small},
    source node/.style={flow node, draw=setSborder, fill=setSsource, text=setStext, line width=1.1pt},
    sink node/.style={flow node, draw=setTborder, fill=setTsink, text=setTtext, line width=1.1pt},
    s node internal/.style={flow node, draw=setSborder!80, fill=setSnode, text=setStext, line width=0.9pt},
    t node internal/.style={flow node, draw=setTborder!80, fill=setTnode, text=setTtext, line width=0.9pt},
    trans node/.style={flow node, fill=yellow!15},
    flow edge/.style={->, thick},
    sat edge/.style={->, line width=1.5pt, cutviolet},
    res edge/.style={->, blue, dashed},
    rev edge/.style={->, red, dashed, line width=1.1pt},
    cut S/.style={dashed, draw=setSborder, fill=setSbg, rounded corners, thick},
    cut T/.style={dashed, draw=setTborder, fill=setTbg, rounded corners, thick},
    % ---------------------------------------------------------------------
    % Rótulos de Partição e Legendas
    % ---------------------------------------------------------------------
    cut label S/.style={font=\footnotesize\bfseries, text=setStext, anchor=north west},
    cut label T/.style={font=\footnotesize\bfseries, text=setTtext, anchor=north east},
    bip label L/.style={font=\large\bfseries, blue!80!black},
    bip label R/.style={font=\large\bfseries, orange!80!black},
    % ---------------------------------------------------------------------
    % Network Simplex e Árvores de Base
    % ---------------------------------------------------------------------
    tree edge/.style={->, line width=1.8pt, blue!80!black},
    nontree edge/.style={->, gray, dashed},
    pivot enter/.style={->, red, dashed, line width=1.5pt},
    pivot leave/.style={->, gray, dashed, line width=1.2pt},
    % ---------------------------------------------------------------------
    % Grafos em Camadas e Distâncias
    % ---------------------------------------------------------------------
    active node/.style={flow node, fill=yellow!30},
    gap node/.style={flow node, fill=orange!40},
    layer line/.style={gray!40, dashed},
    % ---------------------------------------------------------------------
    % Grafos Bipartidos e Emparelhamento
    % ---------------------------------------------------------------------
    bip L/.style={circle, draw, fill=blue!15, minimum size=7.5mm, inner sep=1pt, font=\small},
    bip R/.style={circle, draw, fill=orange!15, minimum size=7.5mm, inner sep=1pt, font=\small},
    cover node/.style={double, double distance=1.2pt},
    match edge/.style={-, line width=1.8pt, blue!80!black},
    unmatch edge/.style={-, gray},
    % ---------------------------------------------------------------------
    % Segmentação de Imagens e Caixas Conceituais
    % ---------------------------------------------------------------------
    pixel node/.style={circle, draw, fill=white, minimum size=8mm, font=\footnotesize},
    nlink/.style={-, thick},
    tlink s/.style={->, green!60!black, dashed},
    tlink t/.style={->, red!70!black, dashed},
    concept box/.style={draw, rounded corners=4pt, minimum height=1.4cm, align=center, font=\small},
    legend node/.style={circle, draw, minimum size=4mm, inner sep=0pt},
    % ---------------------------------------------------------------------
    % Rótulos Padronizados para Arestas
    % ---------------------------------------------------------------------
    edge lbl/.style={font=\footnotesize, fill=white, inner sep=1.2pt, pos=0.5},
    edge lbl sloped/.style={edge lbl, sloped},
    edge lbl above/.style={edge lbl sloped, above=2.8pt},
    edge lbl below/.style={edge lbl sloped, below=2.8pt},
    edge lbl right/.style={edge lbl, right=3pt},
    edge lbl left/.style={edge lbl, left=3pt},
    edge lbl above norm/.style={edge lbl, above=3pt},
    edge lbl below norm/.style={edge lbl, below=3pt},
    edge lbl green/.style={font=\footnotesize, fill=green!10, inner sep=1.2pt, pos=0.5},
    edge lbl red/.style={font=\footnotesize, fill=red!10, inner sep=1.2pt, pos=0.5},
    stage box/.style={draw=blue!80!black, fill=blue!5, rounded corners=4pt, line width=1pt, align=center, font=\footnotesize, minimum height=1.4cm, text width=2.4cm},
    main/.style={flow node},
    % ---------------------------------------------------------------------
    % Pílulas da Capa e Badges
    % ---------------------------------------------------------------------
    cover badge/.style={
        draw,
        rounded corners=3pt,
        fill=topbg!95,
        font=\fontsize{7.5}{9.5}\selectfont\bfseries,
        minimum height=0.46cm,
        text height=5.0pt,
        text depth=0.5pt,
        inner xsep=5pt,
        inner ysep=0pt
    }
}

% =========================================================================
% OPERADORES MATEMÁTICOS E MACROS AUXILIARES
% =========================================================================
\DeclareMathOperator{\ccap}{cap}                  % [OPERADOR] Capacidade de corte funcional: cap(S,T)
\DeclareMathOperator{\cost}{custo}                % [OPERADOR] Custo total de um fluxo: custo(f)
\DeclareMathOperator{\dist}{dist}                 % [OPERADOR] Distância em caminhos mínimos: dist(u,v)

\newcommand{\cred}[2][]{\bar{w}_{#1}(#2)}         % [NOTACAO] Custo reduzido: \cred[\pi]{u,v}
\newcommand{\cres}[2][f]{c_{#1}(#2)}              % [NOTACAO] Capacidade residual: \cres{u,v}
\newcommand{\Dres}[1][f]{D_{#1}}                  % [NOTACAO] Digrafo residual: \Dres
\newcommand{\code}[1]{\lstinline{#1}}             % [MACRO] Código inline ergonômico
\newcommand{\bannertext}[1]{\textbf{\textsf{#1}}} % [MACRO] Cabeçalho destacado de tabela
